\input{../tex-inputs/latex-leseansicht-vorspann}

\section[Gustav Schwarzkopf, Malvine und Vincenz Chiavacci an Arthur Schnitzler, 20. 2. 1899]{L04298 Gustav Schwarzkopf, Malvine und Vincenz Chiavacci an Arthur Schnitzler, 20. 2. 1899}
\nopagebreak\mylabel{L04298v}
\rehead{ }\normalsize\beginnumbering\briefempfaengerindex{Schnitzler, Arthur@\textsc{Schnitzler, Arthur}!zzzChiavacci, Vincenz@\emph{von Vincenz Chiavacci}!1899-02-201@{20. 2. 1899}|(be}\briefempfaengerindex{Schnitzler, Arthur@\textsc{Schnitzler, Arthur}!zzzSchönherr, Malvine@\emph{von Malvine Schönherr}!1899-02-201@{20. 2. 1899}|(be}\briefempfaengerindex{Schnitzler, Arthur@\textsc{Schnitzler, Arthur}!zzzSchwarzkopf, Gustav@\emph{von Gustav Schwarzkopf}!1899-02-201@{20. 2. 1899}|(be}
\toendnotes[C]{\smallbreak\pagebreak[2]}
\correspDesc{Versand  durch Gustav Schwarzkopf, Malvine Chiavacci, Vincenz Chiavacci am 20. 2. 1899 in München
\newline{}Übermittlung  am 20. 2. 1899 in München
\newline{}Zustellung  am 21. 2. 1899 in Wien
\newline{}Erhalt  durch Arthur Schnitzler am 21. 2. 1899 in Wien}\toendnotes[C]{\smallbreak}
\Standort{CUL, Schnitzler, B 96a.}
\physDesc{Bildpostkarte, 135 Zeichen
\newline{}Handschrift Gustav Schwarzkopf: Bleistift, deutsche Kurrent
\newline{}Handschrift Malvine Schönherr: Bleistift, deutsche Kurrent
\newline{}Handschrift Vincenz Chiavacci: Bleistift, lateinische Kurrent
\newline{}Versand: 1) Stempel: »\nobreak{}\oindex{München@\textbf{München}|pwk}Muenchen 1, 20 Feb 99, 12–1Nm\nobreak{}«.   2) Stempel: »\nobreak{}\oindex{Wien@\textbf{Wien}, \emph{Verwaltungsgebiet}|pwk}Wie{[}n{]}, 21. 2. {[}99{]}, 9.V, Bes{[}tellt{]}\nobreak{}«. }\pstart{}{\pb}Herrn \textsc{Dr Arthur
                     Schnitzler}\pend{}\pstart{}\textsc{Wien\oindex{Wien@\textbf{Wien}, \emph{Verwaltungsgebiet}|pw}}\pend{}\pstart{}IX Frankgaſſe 1\oindex{Wien@\textbf{Wien}!IX., Alsergrund@\textbf{IX., Alsergrund}!Frankgasse 1@\textbf{Frankgasse 1}, \emph{Wohngebäude}|pw}.\pend{}{\bigskip}
\pstart
           \noindent{}\centering{}{\pb}\textcolor{gray}{\textbf{Gruss aus dem Nürnberger-Bratwurstglöckl\oindex{Nürnberger Bratwurst Glöckl am Dom@\textbf{Nürnberger Bratwurst Glöckl am Dom}, \emph{Restaurant}|pw}}}\pend
           
\pstart
           \textcolor{gray}{\textbf{München, Frauenplatz 9\oindex{Frauenplatz 9@\textbf{Frauenplatz 9}, \emph{Gebäude}|pw}.}}\pend
           \vspace{1em}
\pstart
           \raggedleft{}{\pb}20/2 1899\pend
           \vspace{0.5em}
\pstart
           Herzlichſt{\\[\baselineskip]}\spacefill\mbox{Gustav Schwzk}\pend
           \leftskip=0em{}\selectlanguage{ngerman}\vspace{1em}
\pstart
           {[}hs. Schönherr:{]} Herzlichen Gruß{\\[\baselineskip]}\spacefill\mbox{Malwine Chiavacci}\pend
           \leftskip=0em{}\selectlanguage{ngerman}\vspace{1em}\pstart {[}hs. Chiavacci:{]} Prosit Blume!, \spacefill\mbox{V. Chiavacci}\pend{}\selectlanguage{ngerman}\endnumbering\briefempfaengerindex{Schnitzler, Arthur@\textsc{Schnitzler, Arthur}!zzzChiavacci, Vincenz@\emph{von Vincenz Chiavacci}!1899-02-201@{20. 2. 1899}|)be}\briefempfaengerindex{Schnitzler, Arthur@\textsc{Schnitzler, Arthur}!zzzSchönherr, Malvine@\emph{von Malvine Schönherr}!1899-02-201@{20. 2. 1899}|)be}\briefempfaengerindex{Schnitzler, Arthur@\textsc{Schnitzler, Arthur}!zzzSchwarzkopf, Gustav@\emph{von Gustav Schwarzkopf}!1899-02-201@{20. 2. 1899}|)be}\mylabel{L04298h}
\begin{anhang}
\end{anhang}\newcommand{\dateiname}{L04298}\newcommand{\titel}{Gustav Schwarzkopf, Malvine und Vincenz Chiavacci an Arthur Schnitzler, 20. 2. 1899}\newcommand{\editorInnen}{Herausgegeben von Jahnke, SelmaMüller, Martin Anton}\input{../tex-inputs/latex-leseansicht-abspann}
